%! AP Statistics — Unit 4, Lesson 4.1 — Homework
\documentclass[10pt]{article}
\usepackage{apstats-article}
\usepackage{apstats-boxes}
\usepackage{grading}
\gradingassignment{apstats-u04-l01-hw}
\gradingcourse{advanced_statistics}

\begin{document}

\pageheader{Unit 4, Lesson 4.1}{Homework}

\vspace{0.1in}

\begin{practicebox}
\small
For each pattern below, (a) write a statistical question it suggests and (b) explain
whether chance alone is a plausible explanation.

\vspace{8pt}
\textbf{1.} A professional tennis player wins 10 consecutive matches on clay courts.
Her overall win rate on clay is 65\%.

\begin{enumerate}[label=(\alph*), leftmargin=2em, itemsep=8pt]
  \item Statistical question: \gradedhide{\writeline\writeline}\gradedright{\raisebox{-10pt}{\gradedblank[id=hw1a, type=rubric, points=2, display={Win rate above 65\%, or chance?}]{2.3in}{0.22in}}}
  \item Could chance alone explain this streak? \gradedhide{\writeline\writeline}\gradedright{\raisebox{-10pt}{\gradedblank[id=hw1b, type=rubric, points=2, display={Yes: $(0.65)^{10} \approx 0.013$, rare but possible.}]{3.5in}{0.22in}}}
\end{enumerate}

\vspace{8pt}
\textbf{2.} A genetics researcher notices that in a sample of 50 plants, 34 have
purple flowers and 16 have white flowers. The theoretical ratio predicted by Mendel's
laws is 3:1 (purple to white).

\begin{enumerate}[label=(\alph*), leftmargin=2em, itemsep=8pt]
  \item What is the observed ratio? \makebox[3cm][r]{\smash{\gradedblank[id=hw2a1, type=short, points=1, answer={34:16 = 2.125:1}, check=text]{1.9cm}{0.15in}}} \quad
        What ratio does Mendelian genetics predict? \makebox[3cm][r]{\smash{\gradedblank[id=hw2a2, type=short, points=1, answer=3:1, check=text]{1.9cm}{0.15in}}}
  \item Statistical question: \gradedhide{\writeline\writeline}\gradedright{\raisebox{-10pt}{\gradedblank[id=hw2b, type=rubric, points=2, display={Is the true ratio 3:1?}]{2.4in}{0.22in}}}
  \item Is the difference between the observed and expected counts large enough to doubt
        the 3:1 ratio? Explain your reasoning. \gradedhide{\writeline\writeline}\gradedright{\raisebox{-10pt}{\gradedblank[id=hw2c, type=rubric, points=2, display={No: chance explains it.}]{1.7in}{0.22in}}}
\end{enumerate}

\vspace{8pt}
\textbf{3.} A weather report shows that the last 6 Sundays in a city have all been
sunny. The historical probability of a sunny Sunday in that city is 0.45.

\begin{enumerate}[label=(\alph*), leftmargin=2em, itemsep=8pt]
  \item Statistical question: \gradedhide{\writeline\writeline}\gradedright{\raisebox{-10pt}{\gradedblank[id=hw3a, type=rubric, points=2, display={Is $P(\text{sunny Sunday}) > 0.45$, or chance?}]{4.0in}{0.22in}}}
  \item If Sundays were independent, what is the approximate probability of 6 sunny
        Sundays in a row? (Hint: multiply the probabilities.) Show your work.\gradedbelow[r]{0.02in}{\gradedblank[id=hw3b, type=short, points=2, answer=0.0083, check=numeric, tolerance=0.0003]{0.9in}{0.24in}}
        \vspace{20pt}
  \item Does this make the streak surprising? \gradedhide{\writeline}\gradedright{\raisebox{-10pt}{\gradedblank[id=hw3c, type=rubric, points=2, display={Yes: only about a 0.8\% chance.}]{3.6in}{0.22in}}}
\end{enumerate}

\vspace{8pt}
\textbf{4.} (AP-Style) A coin is flipped 40 times and lands heads 26 times. Which
of the following best explains why this does NOT necessarily mean the coin is unfair?

\begin{enumerate}[label=(\Alph*), leftmargin=2em, itemsep=4pt]
  \item 26 heads in 40 flips is impossible with a fair coin.
  \item Random variation can produce results that differ from expected values, even with a
        fair coin.
  \item The coin must be flipped at least 100 times before conclusions can be drawn.
  \item A heads proportion of 0.65 is within 0.01 of 0.50, so it is close enough.
        \hfill\gradedright{\textit{Letter:}~\gradedblank[id=hw4, type=choice, points=1, answer=B]{0.4in}{0.2in}}
\end{enumerate}

\vspace{8pt}
\textbf{5.} Reflection: Why is it dangerous to assume every pattern in data has a
non-random cause? Give one real-world example where this mistake could have harmful
consequences.\\[6pt]
\gradedhide{\gradedhere{\kern2.1in\gradedblank[id=hw5, type=rubric, points=2, display={False causes mislead; e.g.\ a chance disease cluster blamed on a vaccine.}]{\linewidth+6pt-2.1in}{\gradedlinesheight{1}+6pt}}\writeline\writeline\writeline}

\vspace{8pt}
\textbf{Extension / Preview:} Find one news headline that describes a ``surprising''
streak, cluster, or trend. Write the headline below, then write a statistical question
it suggests. Be ready to share in class.\\[6pt]
\writeline\writeline\writeline
\end{practicebox}

\end{document}
